\chapter{Periodic Line, Fourier Duality, and Polarization}

We begin not with the infinite line, but with a particle on a line whose ends are identified. That is the deliberate opening move of the lecture. One first solves a tame problem in which momentum is discretized and normalization is honest, and only then lets the size of the box go to infinity. From there the lecture unfolds in a very characteristic sequence: orthogonality, completeness, Fourier transforms, two equivalent representations of the same state, and then the incompatibility of position and momentum. After the break the same logic is repeated in miniature for the polarization of a photon. These notes follow Leonard Susskind's progression closely, with companion-note curation indebted to LazyingArt LLC and Video2Book.

\section{A Particle on a Periodic Line}

A periodic line can be pictured in two ways. One may imagine a particle moving on a literal circle, but one may also imagine an ordinary line segment together with the rule that when the particle exits one end, it reappears at the other. The second picture is the one the lecture emphasizes: the geometry itself is not sacred; what matters is the identification rule. The infinite line will later be recovered as the limit in which the length grows without bound.

The quantum-mechanical translation of that rule is immediate. If the line has length \(L\), then the wave function must satisfy
\begin{equation}
\psi(x)=\psi(x+L).
\end{equation}
This is the first mathematical statement of the lecture.

\begin{figure}[htbp]
\centering
\includegraphics[width=0.72\textwidth]{lecture_05_figure_02.png}
\caption{Periodicity condition on a line of length \(L\).}
\end{figure}

The board here is useful evidence: the equation is written together with a finite line segment labeled \(L\), so the notation and the geometry arrive as one package. It is worth keeping that pairing visible.

\begin{figure}[htbp]
\centering
\begin{tikzpicture}[scale=1]
\node at (3,1.15) {\(\psi(x)=\psi(x+L)\)};
\draw[thick] (0,0) -- (6,0);
\node at (3,-0.35) {\(L\)};
\end{tikzpicture}
\caption{Cleaned redraw of the periodic-line schematic.}
\end{figure}

At this point the lecture backs up one step and says, in effect: let us write
\begin{equation}
L=2\pi r.
\end{equation}
This is introduced purely for notational convenience. One may think of the periodic line as a circle of circumference \(2\pi r\), but the point is not to insist that the particle is really moving on a circle. The point is to have a compact notation that makes the periodicity of plane waves easy to read.

Now the periodicity condition has implications for momentum. The wave function for a momentum eigenstate is first written in the naive form \(e^{ipx}\), and then immediately corrected:
\begin{equation}
e^{ipx/\hbar}=e^{ikx},
\qquad
k=\frac{p}{\hbar}.
\end{equation}
Susskind then switches to \(k\)-notation and stays with it, except for occasional lapses when \(\hbar\) is tacitly set to \(1\).

The plane wave, however, is not yet normalized. On the interval \(0\le x\le 2\pi r\), the normalized momentum eigenfunction is
\begin{equation}
\psi_k(x)=\frac{e^{ipx/\hbar}}{\sqrt{2\pi r}}
=\frac{e^{ikx}}{\sqrt{2\pi r}}.
\end{equation}

\begin{figure}[htbp]
\centering
\includegraphics[width=0.62\textwidth]{lecture_05_figure_03.png}
\caption{Normalized plane wave in \(p\)- and \(k\)-notation.}
\end{figure}

The frame is useful because it records the final board form: the normalized expression is shown side by side in \(p/\hbar\)- and \(k\)-notation, precisely at the moment when the lecture is explaining why the square root appears in the denominator.

\subsection{Question \& Answer}

Why is there a square root in the denominator? Because the normalization condition is imposed on \(\psi^*\psi\), not on \(\psi\) itself. For the bare plane wave \(e^{ikx}\) we have
\begin{equation}
|e^{ikx}|^2=1,
\end{equation}
and so over one full period,
\begin{equation}
\int_0^{2\pi r} |e^{ikx}|^2\,dx = 2\pi r.
\end{equation}
If we divide the wave function by \(\sqrt{2\pi r}\), then the square of that denominator is exactly \(2\pi r\), and the normalization condition becomes
\begin{equation}
\int_0^{2\pi r} |\psi_k(x)|^2\,dx = 1.
\end{equation}
That is the whole reason for the square root. As Susskind remarks in class, the factor is not mysterious; it is simply where all those later \(2\pi\)'s come from.

\section{Allowed Momenta and Discrete Orthogonality}

Once the normalized plane wave has been written down, the periodicity condition can be imposed on the eigenfunction itself. This is where the periodic box begins to earn its keep. We require
\begin{equation}
e^{ik(x+2\pi r)}=e^{ikx}.
\end{equation}
Dividing out the common factor \(e^{ikx}\), we obtain
\begin{equation}
e^{i2\pi rk}=1,
\end{equation}
and therefore the phase accumulated in one trip around the circle must be an integer multiple of \(2\pi\):
\begin{equation}
2\pi r\,k = 2\pi n,
\qquad n\in \mathbb{Z}.
\end{equation}
Thus
\begin{equation}
rk=n,
\qquad
k=\frac{n}{r}.
\end{equation}
The allowed values of \(k\) are therefore discrete, with spacing
\begin{equation}
\Delta k=\frac{1}{r}.
\end{equation}

This is the staging ground for the continuum. At finite \(r\) the momentum spectrum is a lattice of evenly spaced values. As \(r\to\infty\), the spacing goes to zero, the lattice becomes dense, and we are led back toward the infinite line.

Before taking that limit, however, the lecture pauses and says: before we do that, notice something first. These functions are eigenfunctions of a Hermitian operator, and therefore different eigenvalues must correspond to orthogonal eigenfunctions.

\begin{theorem}
If two eigenvectors of a Hermitian operator have distinct eigenvalues, then they are orthogonal.
\end{theorem}

The lecture then checks the statement explicitly for the periodic plane waves. For allowed values \(k=n/r\) and \(k'=n'/r\),
\begin{equation}
\frac{1}{2\pi r}\int_0^{2\pi r} e^{ikx}e^{-ik'x}\,dx
=
\delta_{nn'}.
\end{equation}
If \(k=k'\), the exponentials cancel and the integrand is \(1\), so the integral gives \(1\). If \(k\neq k'\), the product oscillates, with positive contributions canceling negative ones over a full period, and the integral gives \(0\).

This is exactly the point where the lecture stops being merely about a list of allowed momenta. These normalized eigenfunctions are already an orthonormal set; in other words, they are already behaving like a basis of vectors.

\subsection{Question \& Answer}

Is the orthogonality relation a definition? No. This is one of the places where the lecture deliberately slows down. The integral is something we calculate. The deeper statement is the theorem: eigenvectors of a Hermitian operator with different eigenvalues are orthogonal. One may view the explicit integral as a confirmation, in this example, of the general theorem. But the theorem is prior.

\section{From the Box to the Continuum}

Now the lecture takes the next step: the transition from the finite periodic box to the continuum. This is also the point where the spoken derivation becomes a little tangled, so it is worth writing the clean version carefully.

We begin with the box-normalized states
\begin{equation}
\psi_k(x)=\frac{e^{ikx}}{\sqrt{2\pi r}}.
\end{equation}
At finite \(r\), these are normalized to \(1\), and their inner products give Kronecker deltas on the integer labels \(n\).

To prepare the continuum limit, Susskind rescales the states by \(\sqrt r\). Define
\begin{equation}
\phi_k(x)=\sqrt r\,\psi_k(x)=\frac{e^{ikx}}{\sqrt{2\pi}}.
\end{equation}
These are no longer normalized to \(1\) on the finite circle. That is intentional. They are being prepared for a different normalization, one suitable to a continuous spectrum.

What does this rescaling do? For \(k\neq k'\), nothing essential changes: the inner product is still \(0\). But for \(k=k'\), the inner product, which used to be \(1\), becomes \(r\). As \(r\) becomes very large, the equal-\(k\) spike grows while the spacing \(\Delta k=1/r\) shrinks. This is the mechanism by which the Kronecker delta becomes a Dirac delta. In the continuum notation we write
\begin{equation}
\langle k|k'\rangle=\delta(k-k').
\end{equation}

At the same time, the momentum eigenfunction in the position basis becomes
\begin{equation}
\langle x|k\rangle=\frac{e^{ikx}}{\sqrt{2\pi}},
\qquad
\langle k|x\rangle=\frac{e^{-ikx}}{\sqrt{2\pi}}.
\end{equation}
The second equation is the complex conjugate of the first. The notation is now adapted to a continuum basis, and the \(r\) has disappeared from the kernel.

The lecture does not attempt a rigorous theory of distributions here. Its real point is more practical: once the rescaling has been performed, one may pass comfortably to the continuum notation in which sums become integrals and Kronecker deltas become delta functions.

\section{Dyads, Completeness, and Fourier's Trick}

Before going on to Fourier transforms, the lecture inserts a simple but indispensable algebraic fact. The point is to make the coming manipulations lawful. We are not going to pull Fourier transforms out of thin air; we are going to insert the identity operator in a useful way.

The inner product \(\langle a|b\rangle\) is a number. The outer product
\begin{equation}
|b\rangle\langle a|
\end{equation}
is an operator. Its action on a vector \(|c\rangle\) is
\begin{equation}
\bigl(|b\rangle\langle a|\bigr)|c\rangle
=
|b\rangle\,\langle a|c\rangle.
\end{equation}
The bra \(\langle a|\) first contracts with \(|c\rangle\) to produce a number, and that number multiplies the vector \(|b\rangle\). This is the dyad.

Now let \(\{|n\rangle\}\) be an orthonormal basis. Any vector may be expanded as
\begin{equation}
|v\rangle=\sum_n v_n |n\rangle,
\qquad
v_n=\langle n|v\rangle.
\end{equation}
Substituting the second formula into the first gives
\begin{equation}
|v\rangle=\sum_n |n\rangle\langle n|v\rangle.
\end{equation}
Since this holds for every vector \(|v\rangle\), the operator in front must be the identity:
\begin{equation}
I=\sum_n |n\rangle\langle n|.
\end{equation}

\begin{figure}[htbp]
\centering
\includegraphics[width=0.52\textwidth]{lecture_05_figure_04.png}
\caption{Resolution of the identity in a discrete basis.}
\end{figure}

This is exactly the centered board formula captured in the retained frame. In the continuum, the same statement becomes
\begin{equation}
I=\int dk\,|k\rangle\langle k|,
\qquad
I=\int dx\,|x\rangle\langle x|.
\end{equation}
The notation changes, but the idea does not. A complete basis resolves the identity.

This is the trick the lecture now uses over and over. If a calculation looks hard in one basis, insert the identity written in another basis.

\section{Position and Momentum as Two Representations}

Let \(|\psi\rangle\) be the state of a particle on a line. The position-space wave function is defined by
\begin{equation}
\psi(x)=\langle x|\psi\rangle.
\end{equation}
The momentum-space wave function is defined by
\begin{equation}
\tilde\psi(k)=\langle k|\psi\rangle.
\end{equation}
These are not two different states. They are two different descriptions of the same state vector.

The first key observation is that the wave function of a definite momentum state \(|k\rangle\) in the position basis is just
\begin{equation}
\langle x|k\rangle=\frac{e^{ikx}}{\sqrt{2\pi}}.
\end{equation}
This is already a strong hint that position and momentum are going to be related by Fourier theory.

To compute \(\tilde\psi(k)\) from \(\psi(x)\), insert the identity in the position basis:
\begin{align}
\tilde\psi(k)
&=\langle k|\psi\rangle \\
&=\int dx\,\langle k|x\rangle\langle x|\psi\rangle \\
&=\frac{1}{\sqrt{2\pi}}\int dx\,e^{-ikx}\psi(x).
\end{align}
That is the Fourier transform. The minus sign in the exponential appears because \(\langle k|x\rangle\) is the complex conjugate of \(\langle x|k\rangle\).

Now go in the other direction. Insert the identity in the momentum basis:
\begin{align}
\psi(x)
&=\langle x|\psi\rangle \\
&=\int dk\,\langle x|k\rangle\langle k|\psi\rangle \\
&=\frac{1}{\sqrt{2\pi}}\int dk\,e^{ikx}\tilde\psi(k).
\end{align}
This is the inverse Fourier transform. The asymmetry between plus and minus signs is the familiar asymmetry of the Fourier-transform pair.

These formulas are, of course, older than quantum mechanics. The lecture explicitly notes that. But quantum mechanics gives them a new interpretation: they relate two descriptions of one and the same physical state. If somebody hands us \(\psi(x)\), we can immediately compute \(\tilde\psi(k)\); if somebody hands us \(\tilde\psi(k)\), we can recover \(\psi(x)\). The information is the same.

It is also worth recording the operator actions in the two representations. In the position representation,
\begin{align}
\hat X\,\psi(x) &= x\,\psi(x), \\
\hat k\,\psi(x) &= -i\frac{\partial}{\partial x}\psi(x).
\end{align}
In the momentum representation,
\begin{align}
\hat k\,\tilde\psi(k) &= k\,\tilde\psi(k), \\
\hat X\,\tilde\psi(k) &= i\frac{\partial}{\partial k}\tilde\psi(k).
\end{align}
The board equation for \(\hat k\,\psi(x)\) is partly occluded in the retained frame, but the transcript makes the intended form unmistakable: \(-i\partial_x\) is the operator being identified with momentum in the \(x\)-representation.

\begin{figure}[htbp]
\centering
\includegraphics[width=0.72\textwidth]{lecture_05_figure_06.png}
\caption{Operator actions in position and momentum space.}
\end{figure}

The top line of the board, mapping the abstract state to a wave function, is only partly visible, but the intended pattern is clear enough: the lecture is laying the \(x\)-space and \(k\)-space rules side by side so that the symmetry can be seen rather than merely asserted.

Two remarks from the lecture are worth preserving here. First, \(k\) can certainly be negative; that simply corresponds to momentum in the opposite direction. Second, the symmetry between \(x\) and \(k\) is not something classical mechanics makes obvious at first sight. It looks strange until one has learned to think in terms of conjugate quantities.

\subsection{Question \& Answer}

Why do we represent the state as \(\psi(x)\) or as \(\tilde\psi(k)\), but not as a joint function of \(x\) and \(k\)? Because the two variables correspond to incompatible observables. No vector is simultaneously an eigenvector of \(\hat X\) and \(\hat k\). A state of definite momentum is completely spread out in position, and by the same symmetry a state of definite position is completely spread out in momentum. One may pass back and forth between the two representations, but one does not assign a single wave function of \(x\) and \(k\) as though both quantities were simultaneously sharp. This is one of the places where quantum mechanics replaces ``and'' by ``or.''

\section{Wave Packets and the Deferred Classical Limit}

At this point the lecture pauses and refuses to claim too much. We have not yet established a full bridge to classical mechanics. We have only seen that the quantum quantities called position and momentum are represented by the operators \(\hat X\) and \(\hat k\), and that they are incompatible. So why should they deserve the classical names at all?

Susskind's answer is a promise rather than a proof. The bridge is provided by wave packets. A wave packet is a state whose position-space wave function is concentrated in some finite region, and whose momentum-space wave function is also concentrated, though of course not infinitely sharply, around some finite range of \(k\)-values.

A convenient schematic form is
\begin{equation}
\psi(x)\sim e^{ik_0x}f(x),
\end{equation}
where \(f(x)\) is a localized envelope. The factor \(f(x)\) keeps the state concentrated in position space, while the oscillatory factor \(e^{ik_0x}\) makes it look locally like a plane wave, and hence gives it an approximate momentum \(k_0\).

The lecture stresses three points. First, the envelope need not be exactly Gaussian, though a Gaussian is the standard example. Second, the packet is as sharp in both position and momentum as the uncertainty principle allows, but never sharper. Third, for sufficiently smooth potentials, such packets tend to propagate in a way that resembles classical motion: the center of the packet in position space and the center of the packet in momentum space move approximately as a classical particle would move.

But all of that is postponed as a theorem to be proved later. Here the packet appears only as a preview, the first quantum object that looks even vaguely classical.

\section{Photon Polarization: The Same Logic in Two Dimensions}

After the break, the lecture makes a brief return to operator language. If \(A\) and \(B\) are complex vectors, then the inner product can be complex, and the outer product \(|A\rangle\langle B|\) need not be Hermitian. One should think in terms of Hermitian conjugation, not merely ordinary complex conjugation. This remark is short and somewhat noisy in the transcript, but its function is clear: it reopens the operator formalism just before the lecture pivots to a simpler system.

That simpler system is the polarization of a photon. The physical setup is concrete. A polarizer transmits one polarization and blocks, reflects, or absorbs the orthogonal one. For a single photon the result is discrete: the photon either gets through or it does not. The apparatus may therefore be viewed in two ways. It prepares a state, because the transmitted photon emerges in a definite polarization. And it measures a property, because the next polarizer can test what that polarization is.

In the lecture the classical electromagnetic-wave picture is used only long enough to define what is meant by polarization. The electric field oscillates in some transverse direction, and that direction labels the polarization. But once the single-photon description begins, the crucial point is that the outcome is indivisible. One photon does not partly pass and partly reflect. It does one or the other.

Now ignore the motion of the photon and focus only on its polarization state. Choose two orthogonal directions, \(x\) and \(y\), and denote the corresponding states by
\begin{equation}
|x\rangle \leftrightarrow
\begin{pmatrix}
1\\
0
\end{pmatrix},
\qquad
|y\rangle \leftrightarrow
\begin{pmatrix}
0\\
1
\end{pmatrix}.
\end{equation}
These are orthonormal basis vectors:
\begin{equation}
\langle x|x\rangle=1,
\qquad
\langle y|y\rangle=1,
\qquad
\langle x|y\rangle=0.
\end{equation}
The orthogonality is physically motivated. A single experiment can distinguish the two configurations, and so the corresponding vectors must be orthogonal.

Now define an observable associated with an \(xy\)-oriented polarizer. Susskind chooses a convenient convention: assign \(+1\) to one outcome and \(-1\) to the other. The numbers themselves are arbitrary; the apparatus and the notebook convention are what matter. In the chapter we follow the lecture's mathematical content and write
\begin{equation}
\hat P_{xy}|x\rangle=+|x\rangle,
\qquad
\hat P_{xy}|y\rangle=-|y\rangle.
\end{equation}
This completely specifies the operator. In matrix form,
\begin{equation}
\hat P_{xy}=
\begin{pmatrix}
1 & 0\\
0 & -1
\end{pmatrix}.
\end{equation}
Equivalently, in the dyad language already introduced earlier,
\begin{equation}
\hat P_{xy}=|x\rangle\langle x|-|y\rangle\langle y|.
\end{equation}
That is a useful formula, because it ties the two-state polarization example back to the earlier discussion of outer products and completeness.

\subsection{Question \& Answer}

What state comes out of a \(45^\circ\) polarizer? The lecture handles this in exactly the right spirit: not by starting with a matrix, but by asking what the state ought to do. A photon emerging from a \(45^\circ\) polarizer is neither purely \(x\)-polarized nor purely \(y\)-polarized. But if we send it into an \(x\)- or \(y\)-aligned polarizer, we expect equal probabilities for the two outcomes. The natural guess is therefore the equal-amplitude superposition
\begin{equation}
|+\rangle=\frac{|x\rangle+|y\rangle}{\sqrt2}.
\end{equation}
Its orthogonal partner is
\begin{equation}
|-\rangle=\frac{|x\rangle-|y\rangle}{\sqrt2}.
\end{equation}
In column-vector form,
\begin{equation}
|+\rangle \leftrightarrow
\begin{pmatrix}
1/\sqrt2\\
1/\sqrt2
\end{pmatrix},
\qquad
|-\rangle \leftrightarrow
\begin{pmatrix}
1/\sqrt2\\
-1/\sqrt2
\end{pmatrix}.
\end{equation}

The lecture then insists on doing the obvious probability calculation in the formalism, just to show that the machinery really works. Suppose the photon is in the state \(|+\rangle\), and we ask for the probability that it passes a \(y\)-polarizer. The amplitude is
\begin{align}
\langle y|+\rangle
&=
\left\langle y\middle|
\frac{|x\rangle+|y\rangle}{\sqrt2}
\right\rangle \\
&=
\frac{1}{\sqrt2}
\bigl(\langle y|x\rangle+\langle y|y\rangle\bigr) \\
&=
\frac{1}{\sqrt2}.
\end{align}
Therefore
\begin{equation}
|\langle y|+\rangle|^2=\frac12.
\end{equation}
Likewise,
\begin{equation}
|\langle x|+\rangle|^2=\frac12,
\qquad
|\langle x|-\rangle|^2=\frac12,
\qquad
|\langle y|-\rangle|^2=\frac12.
\end{equation}
So the diagonal states are exactly the states whose \(x\)- and \(y\)-probabilities are equal.

Now define the observable associated with the \(45^\circ\) basis. Let \(|+\rangle\) have eigenvalue \(+1\) and \(|-\rangle\) have eigenvalue \(-1\):
\begin{equation}
\hat P_{45}|+\rangle=+|+\rangle,
\qquad
\hat P_{45}|-\rangle=-|-\rangle.
\end{equation}
The simplest matrix that does this, in the \((|x\rangle,|y\rangle)\) basis, is
\begin{equation}
\hat P_{45}=
\begin{pmatrix}
0 & 1\\
1 & 0
\end{pmatrix}.
\end{equation}
One checks it immediately:
\begin{equation}
\hat P_{45}
\begin{pmatrix}
1/\sqrt2\\
1/\sqrt2
\end{pmatrix}
=
\begin{pmatrix}
1/\sqrt2\\
1/\sqrt2
\end{pmatrix},
\qquad
\hat P_{45}
\begin{pmatrix}
1/\sqrt2\\
-1/\sqrt2
\end{pmatrix}
=
-
\begin{pmatrix}
1/\sqrt2\\
-1/\sqrt2
\end{pmatrix}.
\end{equation}
Equivalently,
\begin{equation}
\hat P_{45}=|+\rangle\langle +|-|-\rangle\langle -|.
\end{equation}

Now the lecture arrives at the real point. The eigenvectors of \(\hat P_{xy}\) are \(|x\rangle\) and \(|y\rangle\). The eigenvectors of \(\hat P_{45}\) are \(|+\rangle\) and \(|-\rangle\). There are no common eigenvectors. Therefore the two observables are incompatible. This is exactly the same logic as before, but now in a two-dimensional vector space rather than on the infinite line.

\subsection{Question \& Answer}

Why is probability not itself an observable? Because an observable, in the technical sense used in the lecture, is something that gives a definite outcome in a single experiment. A photon sent through a polarizer yields a definite record in the notebook: through or not through, or equivalently \(+1\) or \(-1\). Probability does not appear in a single run. It is inferred only after many repeated runs.

The lecture's late discussion is worth recording in the cleaner dyad language. Define the projectors
\begin{equation}
\Pi_x=|x\rangle\langle x|,
\qquad
\Pi_y=|y\rangle\langle y|.
\end{equation}
Then for a state \(|\psi\rangle\),
\begin{equation}
\langle \psi|\Pi_x|\psi\rangle
=
\langle \psi|x\rangle\langle x|\psi\rangle
=
|\langle x|\psi\rangle|^2,
\end{equation}
and similarly
\begin{equation}
\langle \psi|\Pi_y|\psi\rangle
=
|\langle y|\psi\rangle|^2.
\end{equation}
These expectation values are probabilities. But the output of the single-shot experiment associated with \(\Pi_x\) is \(1\) or \(0\), not the probability. The same distinction is what underlies the lecture's final insistence that an observable is tied to a particular apparatus and to the little number one writes down event by event after the system has interacted with that apparatus.

The lecture closes by sharpening the language further: all observables are Hermitian operators, and all Hermitian operators are observables. But the observable is never the abstract probability itself. The observable is the single-run result associated with a definite measurement setup.

\section{Summary}

The first half of the lecture takes a particle on a periodic line, normalizes the plane waves, quantizes the allowed momenta, and then carefully stretches the finite-box normalization into continuum normalization. With dyads and completeness in hand, the lecture derives the Fourier-transform relation between the position and momentum representations and uses that relation to sharpen the incompatibility of \(\hat X\) and \(\hat k\).

The second half rebuilds the same logic in a much simpler system. A polarizer prepares states and measures outcomes; the orthogonal states \(|x\rangle\) and \(|y\rangle\) form a basis; rotated states such as \(|+\rangle\) and \(|-\rangle\) define new observables; and the absence of common eigenvectors expresses incompatibility once again.

What runs through the whole lecture is the same structural lesson. A state may be described in more than one representation, but incompatible observables do not fit into one simultaneous sharp picture. And an observable is not whatever number we infer from an ensemble after the fact; it is the definite outcome recorded in a single experiment by a particular apparatus.